\documentclass[10pt,a4paper]{amsart}
\usepackage[margin=19mm]{geometry}
\usepackage{amsmath,amssymb,mathtools}
\usepackage[scr=rsfs,cal=euler]{mathalfa}
\usepackage{xfrac}
\usepackage[unicode,hidelinks,bookmarks=false]{hyperref}
\newcommand{\ie}{i.e.\ }
\newcommand{\cf}{cf.\ }
\newcommand{\resp}{respectively\ }
\newcommand{\Subsection}{\subsection}
\providecommand{\nobreakdash}{\nobreak\hbox{-}}
\setlength{\emergencystretch}{2em}
\multlinegap0pt
\usepackage{bm}
\usepackage{bbm}
\usepackage{dsfont}
\usepackage{xspace}
\usepackage{cancel}
\usepackage{mathtools}
\usepackage{amsthm}
\makeatletter
\@ifundefined{thm}{\newtheorem{thm}{Thm}}{}
\@ifundefined{prop}{\newtheorem{prop}[thm]{Proposition}}{}
\makeatother
\def\dTC{{\mathop\mathrm{dTC}\,}}
\title[On distributional topological complexity of groups and manif]{On distributional topological complexity of groups and manifolds}
\author{Alexander Dranishnikov}
\date{Source: arXiv:2509.16759v4 (CC BY 4.0)}
\begin{document}
\maketitle

\noindent\emph{Source and licence.} On distributional topological complexity of groups and manifolds. Version arXiv:2509.16759v4 (CC BY 4.0), distributed under the Creative Commons Attribution 4.0 International licence, as recorded in the arXiv licence metadata for this exact version and shown on the abstract page https://arxiv.org/abs/2509.16759v4. Full rights notice: licenses/arxiv-2509.16759-CC-BY-4.0.txt.\par\smallskip

\noindent\textit{Editorial scope.} Counted unit: the Prop whose source label is \texttt{2p-1} in the article (source0.tex lines 575--580, source label \texttt{2p-1}), delivered together with the article's own complete original proof of it (source0.tex lines 581--623, one \texttt{proof} environment of 43 lines). No mathematics is written, repaired or re-derived: statement and proof are the article's own text line for line. Required context retained uncounted: none. Every definition, display and label of those lines is retained unchanged. Omitted from this package and not redistributed: 1--574, 624--1073. Nothing inside a retained excerpt is omitted or altered: every retained body is delivered exactly as the article's lines are written, so no omission receipt of any kind accompanies this package. Citations: the retained excerpts carry no citation command at all, so no bibliography entry of the article is transcribed and no reference is redistributed. Reference adaptation: none was needed and none was made; every reference inside the delivered unit resolves inside it, so no unresolved reference or citation is produced. Printed slips kept literal: no printed word, symbol, hypothesis or number of the counted statement or of its proof was altered. Layout and class: the article's own class is \texttt{amsart} with packages including amssymb, latexsym, amsmath, amscd, array, graphicx, xy; the wrapper is the lane's pinned \texttt{amsart} editorial wrapper at 10pt on A4, which restates the theorem-like environments the retained text needs and adds the article's own macro definitions that the retained text needs (\texttt{\textbackslash dTC}), with extra wrapper packages (bm, bbm, dsfont, xspace, cancel, mathtools). The delivered numbering is the unit's own. Assets: no retained span contains a figure environment, a graphics-inclusion command, a PDF-page inclusion, a table or a TikZ picture; no image file is copied into this package, the package declares no asset dependency and no image file is redistributed with it. Licence: version arXiv:2509.16759v4 of this article is distributed under the Creative Commons Attribution 4.0 International licence, as recorded in the arXiv licence metadata for this exact version and shown on the abstract page https://arxiv.org/abs/2509.16759v4. A full rights notice is delivered at licenses/arxiv-2509.16759-CC-BY-4.0.txt. Characters: every retained excerpt is pure ASCII, so the delivered engine, XeLaTeX with its default Latin Modern text font, reports no missing character.
\begin{prop}\label{2p-1}
When $r$ is odd,
$\dTC(L^m_r)\le 2r-1$ for all $m$ and  $r$ and

$\dTC(L^m_r)\le r-1$ for even $r$.
\end{prop}

\begin{proof}
Let $r$ be odd.
We denote $G=\mathbb Z_r$.
By definition, the space $L^m_r$ is the orbit space of a $G$-action on $S^m$ by isometries, where $m=2n-1$. 
Let $Gx$ and $Gy$ be two orbits. 
Here we will identify $z\in S^m$ with the Dirac measure $\delta_z$ supported at $z$.
Given $x$ and $y$ in $S^m$, we define a $\mathcal B_{2r}$-path in $S^m$ from $x$ to the uniformly distributed measure $\sum_{g\in G}\frac{1}{r}gy\in\mathcal B_r(Gy)$
on the orbit $Gy$
 by the formula
$$
\phi(t)=\sum_{g\in G}\left(\frac{\beta(x,gy)}{2\pi}R^{\alpha
(x,gy)}(t)+\frac{\alpha(x,gy)}{2\pi}R^{\beta(x,gy)}(t)\right),
$$
where $\alpha(x,y)$ and $\beta(x,y)$ are the angles between vectors $x$ and $y$,
$\alpha(x,y)+\beta(x,y)=2\pi$, and where $R^{\alpha(x,y)}(t)$ and $R^{\beta(x,y)}(t)$ are paths defined by the rotations in the plane $A$ containing  $x$ and $y$, taking $x$ to $y$. These are rotations by angles $\alpha(x,y),\beta(x,y)$ in opposite directions. We may assume that $\alpha\le\beta$.  When $\alpha=\beta=\pi$ the rotations are indistinguishable, which does not cause any problem since there is no ordering of summands in the definition of $\phi$.
 



We note that if $\alpha(x,gy)=\pi$ or $\alpha(x,gy)=0$ for some $g\in G$, then the plane $A$ coincides with the plane containing $Gx$, which also coincides with plane containing $Gy$. The path $\phi$ is still well defined. If $\alpha(x,gy)=\pi$, then $$\phi(t)=\frac{1}{2}R^{\alpha(x,gy)}(t)+\frac{1}{2}R^{\beta(x,gy)}(t),$$ where the rotations by $\pi$ are in opposite directions, and there is no ordering between them. If $\alpha(x,gy)=0$, then $R^{\alpha(x,gy)}(t)$ is the identity and the direction of rotation of $R^{\beta(x,gy)}(t)$
is irrelevant since it has coefficient zero.

It is easy to check that each summand in the definition of $\phi$ depends continuously on $x$ and $y$. This implies that the $\mathcal B_{2r}$-path $\phi$  depends continuously on $(x,y)\in S^m\times S^m$.


We define a $\mathcal B_{2r}$-path $\psi$ from $Gx$ to $Gy$ in $L^m_r$ as the image of $\phi$ under the quotient map $q:S^m\to L^m_r$,
$$
\psi_{Gx,Gy}(t)=\sum_{g\in G}\left(\frac{\beta(x,gy)}{\pi}qR^{\alpha(x,gy)}(t)+\frac{\alpha(x,gy)}{\pi}qR^{\beta(x,gy)}(t)\right)
$$
Since $G$ acts by isometries, we have $R^{\alpha(gx,gy)}(t)=gR^{\alpha(x,y)}(t)$ and
$R^{\beta(gx,gy)}(t)=gR^{\alpha(x,y)}(t)$. Therefore $\psi$ does not depend on the choice of the representative $x\in Gx$ in definition of $\phi$. It is easy to see that the map
$$\Psi:L^m_r\times L^m_r\to\mathcal B_{2r}(P(X))$$ defined by $\Psi(Gx,Gy)$ to $\psi_{Gx,Gy}$ is continuous,
and $\psi_{Gx,Gy}(0)=Gx$ and $ \psi_{Gx,Gy}(1)=Gy$.

Now suppose that $r=2k$.  For any vector $z\in Gx$, its negative $-z\in Gx$. We  consider lines $\ell_x$ and $\ell_{gy}$ through the vectors $x$ and $gy$ respectively, and define $\alpha(\ell_x,\ell_{gy})$ and $\beta(\ell_x,\ell_{gy})$ to be the angels between $\ell_x$ and $\ell_{gy}$. Thus,
$\alpha(\ell_x,\ell_{gy})+\beta(\ell_x,\ell_{gy})=\pi.$ Let $H\subset G$ be a subgroup of index 2. We define a $\mathcal B_{2k}$-path (in $\mathbb RP^{2n-1}$) from the line $\ell_x$ to $\sum_{g\in H}\frac{1}{k}\ell_{gy}$ as
$$
\phi(t)=\sum_{g\in H}\left(\frac{\beta(x,gy)}{\pi}R^{\alpha
(x,gy)}(t)+\frac{\alpha(x,gy)}{\pi}R^{\beta(x,gy)}(t)\right).
$$
This defines a $\mathcal B_{2k}$-path in $\mathbb S^m$ from $\frac{1}{2}x+\frac{1}{2}(-x)$ to $\sum_{g\in G}\frac{1}{2k}gy\in\mathcal B_r(Gy)$.
Then we define $\Psi$ as above by projecting to $L_r^m$.
\end{proof}

\end{document}
